\section{External Matrix Interpretation}\label{sec:matrix-background}
Write $V=\Zp^{2n}$ and $\omega=e^{2\pi i/p}$.  Gross's Weyl
operators satisfy $W(v)W(w)=\omega^{\frac12\mathrm{sform}(v,w)}
W(v+w)$, after matching the order of the position and momentum
coordinates~\cite[\S II.1, Eq.~(4)]{gross2006hudson}.
Every Clifford unitary is, up to phase, a Weyl operator followed by
a symplectic implementer~\cite[Theorem~3]{gross2006hudson}.
For finite fields of odd characteristic those implementers can be
chosen as a \emph{linear} Weil representation $\rho$, with exact
covariance $\rho(S)W(v)\rho(S)^{-1}=W(Sv)$
\cite[Theorem~1.3.2]{gurevich2007geometric}.
Thus, writing $U(e,v,S)=\omega^e W(v)\rho(S)$, one obtains
\[
 U(e,v,S)U(f,w,T)
 =U\bigl(e+f+\tfrac12\mathrm{sform}(v,Sw),\,v+Sw,\,ST\bigr).
\]
This is the Heisenberg--Weil representation of the Jacobi group
\cite[\S 1.3]{gurevich2010bounds}, and its multiplication is exactly
our structural extension's.  For $n\geq1$ it is faithful: conjugation
on Weyl operators determines $S$ and then $v$, and the remaining
scalar determines $e$.  Its scalar subgroup is therefore exactly
$\langle\omega I\rangle$; at width zero the representation is simply
$e\mapsto\omega^e$.
Adjoining $-I$ gives a direct product with $\mathbb Z_2$, since $p$ is
odd, with scalar subgroup $\langle-\omega I\rangle$.
These are the two finite scalar conventions used here; the full
unitary normalizer allows every scalar in $U(1)$.
The representation argument is mathematical background, not an Agda
theorem.  An isomorphism of abstract groups would not identify the named gates:
one must also prove that their interpretations are the specific
$\delta_p$-normalised matrices. This gate-by-gate correspondence is
not proved here. Explicit constructions on symplectic generators appear in
Neuhauser~\cite{neuhauser2002metaplectic}.
